\documentclass{article}
\usepackage[a4paper, total={6in, 8in}]{geometry}
\usepackage{graphicx, amssymb, pifont, amsmath}

\title{Tecnología Digital VI - Guía de Ejercicios 4}
\author{Juan Ignacio Elosegui}
\date{Noviembre 2025}

\begin{document}
    \maketitle
    \section*{Ejercicio 1}
        \noindent El algoritmo de \textit{backpropagation} cumple con las siguientes opciones:
        \begin{itemize}
           \item Es un método de regularización para redes neuronales.
           \item Es un algoritmo para el cálculo del gradiente de una función.
        \end{itemize}
    \section*{Ejercicio 2}
        \subsection*{Inciso a)}
            \noindent PyTorch permite diferenciar funciones por medio de \textbf{diferenciación automática}, con el método \textit{torch.autograd}.
        \subsection*{Inciso b)}
            \noindent PyTorch permite:
            \begin{enumerate}
                \item Paralelizar operaciones en múltiples procesadores de CPU y GPU.
                    \subitem \noindent Se llama paralelismo de datos. Distribuye el trabajo computacional.
            \end{enumerate}
    \section*{Ejercicio 3}
        \noindent Sea la función \[ f(w, x, y, z) = -\left(\max(z, 2x) + 2(wy)\right) \tag{1} \] evaluada en el punto \(w = 3, x = 2, y = 4, z = 1\), y la función \[ g(w, x, y, z) = -3((4xy)+\max(w, z)) \tag{2} \] evaluada en el punto \(w = 3, x = 5, y = 2, z = 4\).
        \subsection*{Inciso 1}
            \begin{center}
                includegraphics[width=0.7\linewidth]{figs/3.1.png}
                includegraphics[width=0.7\linewidth]{figs/3.1b.png}
            \end{center}
        \subsection*{Inciso 2}
            
\end{document}
